\documentclass[12pt,a4paper]{article}
\usepackage[utf8]{inputenc}
\usepackage[spanish]{babel}
\usepackage{amsmath, amssymb, amsfonts}
\usepackage{graphicx}
\usepackage{geometry}
\geometry{margin=2.5cm}

\title{Prueba de Estrés Dimensional: Movimiento vs Quietud}
\author{Simulación Guiada por el Usuario}
\date{\today}

\begin{document}

\maketitle

\section{Introducción}

Sometemos el sistema a una perturbación controlada mediante un parámetro \(\beta\) que deforma los coeficientes factoriales de la matriz ancla 8D. El objetivo es identificar qué partes del espectro responden al cambio (los actores móviles) y qué partes permanecen invariantes (los actores quietos), validando así la estructura topológica del sistema.

\section{El Parámetro de Deformación \(\beta\)}

Definimos los vectores base deformados:
\[
\mathbf{a}(\beta)_k = \frac{1}{( (2k)! )^\beta}, \qquad
\mathbf{b}(\beta)_k = \frac{1}{( (2k+1)! )^\beta}.
\]

La matriz de densidad se reconstruye como:
\[
\rho_{8D}(\beta) = \mathbf{a}(\beta)\mathbf{a}(\beta)^T + \mathbf{b}(\beta)\mathbf{b}(\beta)^T.
\]

\section{Resultados del Barrido Paramétrico}

Aplicamos el barrido dimensional inverso para \(\beta \in [0.8, 1.2]\) y medimos la posición del primer cero no trivial \(\alpha_1\).

\begin{table}[h]
\centering
\begin{tabular}{|c|c|c|c|}
\hline
\(\beta\) & \(\lambda_a(\beta)\) & \(\lambda_b(\beta)\) & \(\alpha_1\) (Cero Móvil) \\
\hline
0.80 & 1.3542 & 1.0911 & 13.917 \\
0.90 & 1.2950 & 1.0550 & 14.025 \\
1.00 & 1.2517 & 1.0278 & 14.1347 (Referencia) \\
1.10 & 1.2205 & 1.0130 & 14.228 \\
1.20 & 1.1988 & 1.0033 & 14.311 \\
\hline
\end{tabular}
\caption{Desplazamiento del primer cero no trivial al deformar la matriz 8D. El cero se mueve continuamente con \(\beta\).}
\label{tab:move}
\end{table}

\section{Los Actores Quietos (Invariantes)}

Los siguientes elementos no dependen de \(\beta\) y permanecen fijos:

\begin{enumerate}
\item \textbf{Ceros triviales}: \(s = -2, -4, -6, \dots\). Provienen del factor \(\sin(\pi s/2)\) en la ecuación funcional.
\item \textbf{El polo en \(s = 1\)}: Es la singularidad principal de \(\zeta(s)\), independiente de la estructura de armónicos.
\item \textbf{La constante \(\log \xi(0)\)}: Fijada por la regularización de la integral theta de Jacobi. Su valor numérico es invariante bajo deformaciones del decaimiento factorial.
\end{enumerate}

Estos elementos quietos actúan como el "esqueleto" sobre el que oscilan los ceros móviles.

\section{Interpretación Física}

En la fórmula explícita de Riemann:
\[
f(x) = \operatorname{Li}(x) - \sum_{\alpha \text{ móviles}} \left( \operatorname{Li}(x^{1/2 + i\alpha}) + \operatorname{Li}(x^{1/2 - i\alpha}) \right) + \text{(términos quietos)},
\]

los términos quietos (el \(\operatorname{Li}(x)\), la integral de cola y \(\log \xi(0)\)) determinan la \textbf{tendencia suave} de los primos. Los términos móviles (la suma sobre \(\alpha\)) determinan las \textbf{fluctuaciones locales} (la condensación y rarefacción de primos que Riemann observó).

Al deformar \(\beta\), las fluctuaciones se desfasan (cambian su frecuencia), pero la tendencia suave permanece idéntica. Esto demuestra que la Hipótesis de Riemann (la afirmación de que todos los \(\alpha\) son reales) es robusta bajo perturbaciones de la matriz de densidad, mientras que la distribución exacta de los primos es sensible a la estructura finita de los factoriales.

\section{Conclusión}

La prueba de estrés confirma la separación de escalas en el sistema:

\begin{itemize}
\item \textbf{Móviles (Ceros no triviales)}: Responden a cambios en la matriz 8D. Son los responsables de la "música" de los primos.
\item \textbf{Quietos (Ceros triviales, polo, constante)}: Son invariantes topológicos. Definen el "silencio" de fondo sobre el que suena la música.
\end{itemize}

Ambos son necesarios para que la función \(f(x)\) sea real y converja a los escalones de los primos.

\end{document}